% TikZ 图：矢量加法示意（东3北4得5）
% Day 01：初高中物理的区别与学习方法
% 使用：% TikZ 图：矢量加法示意（东3北4得5）
% Day 01：初高中物理的区别与学习方法
% 使用：% TikZ 图：矢量加法示意（东3北4得5）
% Day 01：初高中物理的区别与学习方法
% 使用：% TikZ 图：矢量加法示意（东3北4得5）
% Day 01：初高中物理的区别与学习方法
% 使用：\input{figures/tikz/vector-addition.tex}
\begin{tikzpicture}[scale=1.4,>=Latex]
  % 画箭头（加长让标注有空间）
  \draw[->,thick] (0,0) -- (3.2,0);
  \draw[->,thick] (3.2,0) -- (3.2,4.2);
  \draw[->,ultra thick,red] (0,0) -- (3.2,4.2);
  % 长度标注放在箭头正中间的上方/右侧，远离坐标轴交汇处
  \node[above,yshift=2pt] at (1.6,0) {$3\ \mathrm{m}$};
  \node[right,xshift=4pt] at (3.2,2.1) {$4\ \mathrm{m}$};
  \node[above left,yshift=4pt,xshift=-2pt] at (1.6,2.1) {$5\ \mathrm{m}$};
  % 方向标注放在箭头末端之外
  \node[below] at (3.2,-0.15) {东};
  \node[right] at (3.5,4.2) {北};
  % 起点
  \node[below left,yshift=-2pt] at (0,0) {起点};
\end{tikzpicture}

\begin{tikzpicture}[scale=1.4,>=Latex]
  % 画箭头（加长让标注有空间）
  \draw[->,thick] (0,0) -- (3.2,0);
  \draw[->,thick] (3.2,0) -- (3.2,4.2);
  \draw[->,ultra thick,red] (0,0) -- (3.2,4.2);
  % 长度标注放在箭头正中间的上方/右侧，远离坐标轴交汇处
  \node[above,yshift=2pt] at (1.6,0) {$3\ \mathrm{m}$};
  \node[right,xshift=4pt] at (3.2,2.1) {$4\ \mathrm{m}$};
  \node[above left,yshift=4pt,xshift=-2pt] at (1.6,2.1) {$5\ \mathrm{m}$};
  % 方向标注放在箭头末端之外
  \node[below] at (3.2,-0.15) {东};
  \node[right] at (3.5,4.2) {北};
  % 起点
  \node[below left,yshift=-2pt] at (0,0) {起点};
\end{tikzpicture}

\begin{tikzpicture}[scale=1.4,>=Latex]
  % 画箭头（加长让标注有空间）
  \draw[->,thick] (0,0) -- (3.2,0);
  \draw[->,thick] (3.2,0) -- (3.2,4.2);
  \draw[->,ultra thick,red] (0,0) -- (3.2,4.2);
  % 长度标注放在箭头正中间的上方/右侧，远离坐标轴交汇处
  \node[above,yshift=2pt] at (1.6,0) {$3\ \mathrm{m}$};
  \node[right,xshift=4pt] at (3.2,2.1) {$4\ \mathrm{m}$};
  \node[above left,yshift=4pt,xshift=-2pt] at (1.6,2.1) {$5\ \mathrm{m}$};
  % 方向标注放在箭头末端之外
  \node[below] at (3.2,-0.15) {东};
  \node[right] at (3.5,4.2) {北};
  % 起点
  \node[below left,yshift=-2pt] at (0,0) {起点};
\end{tikzpicture}

\begin{tikzpicture}[scale=1.4,>=Latex]
  % 画箭头（加长让标注有空间）
  \draw[->,thick] (0,0) -- (3.2,0);
  \draw[->,thick] (3.2,0) -- (3.2,4.2);
  \draw[->,ultra thick,red] (0,0) -- (3.2,4.2);
  % 长度标注放在箭头正中间的上方/右侧，远离坐标轴交汇处
  \node[above,yshift=2pt] at (1.6,0) {$3\ \mathrm{m}$};
  \node[right,xshift=4pt] at (3.2,2.1) {$4\ \mathrm{m}$};
  \node[above left,yshift=4pt,xshift=-2pt] at (1.6,2.1) {$5\ \mathrm{m}$};
  % 方向标注放在箭头末端之外
  \node[below] at (3.2,-0.15) {东};
  \node[right] at (3.5,4.2) {北};
  % 起点
  \node[below left,yshift=-2pt] at (0,0) {起点};
\end{tikzpicture}
